对每个点 $x$，按路径的第一步经过哪个邻点，将其他点分组。每组取一个离 $x$ 最远的点，再连同 $x$ 本身保留距离最大的三项。每项为 $\{d,u,t\}$，分别记录距离、点编号和对应的邻点。距离相等的项也保留；$x$ 本身的 $t=0$，空项的 $u=0$。

第一遍 DFS 先计算子树内的结果。若孩子 $y$ 子树内的最远距离为 $d_y$，则它到 $x$ 的距离为 $d_y+w_{xy}$，对应的 $t=y$。

第二遍 DFS 到达 $x$ 时，已考虑来自父节点方向的最远点。计算孩子 $y$ 的答案时，需要排除 $y$ 子树：若最大项的 $t=y$，就取次大项，否则取最大项。将这项的距离加上边权，加入 $y$ 的结果，并令 $t=x$。

只排除一个邻点并取最大距离时，保留前两大即可。若排除后还要取两个最大距离，例如拼接经过当前点的最长路径，就需要保留前三大。总时间复杂度 $O(n)$。
